\documentclass[xcolor=table,presentation]{beamer}    % ,handout
% colortbl only defines \rowcolor for a single row. xcolor extends this to multiple rows.

\usetheme[sectionoutline]{Aachen} %
\usepackage[ngerman]{babel}
\usepackage[T1]{fontenc}
\usepackage[utf8]{inputenc}


%%%%%%%%%%%%%%%%%%%%%%%%%%%%%%%%%%%%%%%%%%%%%%%%%%%%%%%%%%%%%%%%%%%%%%
% tables
\usepackage{multirow,array,tabularx,rotating}
\usepackage{booktabs}

% math
\usepackage{amsmath,amsthm, amssymb, latexsym, xspace}
\usepackage{bbold}
%\usefonttheme[onlymath]{serif}
%\boldmath

% misc
\usepackage{subfigure}
\usepackage{wasysym}
\usepackage{nameref}

% declare the path(s) where your graphic files are
\graphicspath{./bilder/}
% and their extensions so you won't have to specify these with
% every instance of \includegraphics
\DeclareGraphicsExtensions{.pdf,.jpeg,.png}

% helpers
\newcommand{\argmin}{\operatornamewithlimits{argmin}}
\newcommand{\argmax}{\operatornamewithlimits{argmax}}

%%%%%%%%%%%%%%%%%%%%%%%%%%%%%%%%%%%%%%%%%%%%%%%%%%%%%%%%%%%%%%%%%%%%%%
%%%%%%%%%%%%%%%%%%%%%%%%%%%%%%%%%%%%%%%%%%%%%%%%%%%%%%%%%%%%%%%%%%%%%%

\renewcommand*{\email}{\url{thomas.gatzweiler@rwth-aachen.de}}
% all email address(es) of the authors (used for \TitlePage)

\title[Adaptive Huffman-Kodierung und Anwendungen]{Adaptive Huffman-Kodierung und Anwendugen}

\setbeamertemplate{navigation symbols}{} %disable navigation bar

%% author and in []: shortauthor
\author[Gatzweiler]{Thomas Gatzweiler}
% - Use the \inst{?} command only if the authors have different
%   affiliation.
\institute[RWTH Aachen University] % (optional, but mostly needed)
{
%  \inst{1}%
  \strut Human Language Technology and Pattern Recognition\\
  \strut Computer Science Department, RWTH Aachen University %\\
  %\strut {\tt lehnen@cs.rwth-aachen.de}
}
% - Use the \inst command only if there are several affiliations.
% - Keep it simple, no one is interested in your street address.

\date[12. November 2014]{Proseminar Datenkompression, 12. November 2014, Aachen}

%%%%%%%%%%%%%%%%%%%%%%%%%%%%%%%%%%%%%%%%%%%%%%%%%%%%%%%%%%%%%%%%%%%%%%
% will be set into the PDF document summary
\hypersetup{
  pdftitle={\inserttitle}, 
  pdfauthor={\insertauthor}, 
  bookmarksdepth=subsubsection,  
  % enable automatic page transitions: for endless loop edit in
  % acrobat reader -> preferences -> full screen -> after every X
  % seconds and after last page
  %pdfpageduration = 2, 
  % pdfpagetransition = {Glitter /Di 315 /D 5}  
  % pdfpagetransition = {Box /M /O /D 1},
}

%%%%%%%%%%%%%%%%%%%%%%%%%%%%%%%%%%%%%%%%%%%%%%%%%%%%%%%%%%%%%%%%%%%%%%
%%%%%%%%%%%%%%%%%%%%%%%%%%%%%%%%%%%%%%%%%%%%%%%%%%%%%%%%%%%%%%%%%%%%%%

%%%%%%%%%%%%%%%%%%%%%%%%%%%%%%%%%%%%%%%%%%%%%%%%%%%%%%%%%%%%%%%%%%%%%%
%%%%%%%%%%%%%%%%%%%%%%%%%%%%%%%%%%%%%%%%%%%%%%%%%%%%%%%%%%%%%%%%%%%%%%
\begin{document}
% Literaturverzeichnis wie in der .bib Datei ordnen
\nocite{*}

%%%%%%%%%%%%%%%%%%%%%%%%%%%%%%%%%%%%%%%%%%%%%%%%%%%%%%%%%%%%%%%%%%%%%%
\begin{frame}[label=titlepage]
  \titlepage
\end{frame}

%%%%%%%%%%%%%%%%%%%%%%%%%%%%%%%%%%%%%%%%%%%%%%%%%%%%%%%%%%%%%%%%%%%%%%

% use this, if you use option nosectionoutline
%
% \begin{frame}<*>{\iflanguage{ngerman}{Gliederung}{Outline}}
%       \tableofcontents[subsectionstyle=show/show/show]
% \end{frame}

%%%%%%%%%%%%%%%%%%%%%%%%%%%%%%%%%%%%%%%%%%%%%%%%%%%%%%%%%%%%%%%%%%%%%%

\section{Einführung} \label{sec:intro}
\begin{frame}{\insertsection}
Bla Bla Bla
\end{frame}

%%%%%%%%%%%%%%%%%%%%%%%%%%%%%%%%%%%%%%%%%%%%%%%%%%%%%%%%%%%%%%%%%%%%%%

\begin{frame}[allowframebreaks]{Literatur}
 %\bibliographystyle{alpha}
 \bibliographystyle{i6bibliostyle}
 \bibliography{references}
\end{frame}

%%%%%%%%%%%%%%%%%%%%%%%%%%%%%%%%%%%%%%%%%%%%%%%%%%%%%%%%%%%%%%%%%%%%%%

\section{Huffman-Kodierung}

\subsection{Huffman-Baum}
\begin{frame}[fragile]{\insertsection: \insertsubsection}
Für das Paket bieten wir zwei Optionen an.
 \begin{itemize}
  \item \verb+\usetheme[sectionoutline]{Aachen}+
  \item \verb+\usetheme[nosectionoutline]{Aachen}+
 \end{itemize}
Mit \emph{sectionoutline} werden ab dem zweiten Abschnitt eine Gliederung vor den Abschnitt
eingefügt, wenn das Kommando \verb+\section+ und nicht das Kommando \verb+\section*+ benutzt wurde.
Dies gibt dem Hörer eine Information zum Fortschritt, ordnet das Gehörte ein und sorgt für eine
kleine rhetorische Pause.
\end{frame}

\subsection{Bilder}
\begin{frame}{\insertsection: \insertsubsection}
\begin{center}
 \includegraphics[width=.5\linewidth]{logos/rwth-hks44}
\end{center}
\end{frame}

\subsection{Cite}
\begin{frame}{\insertsection: \insertsubsection}
Wie beschrieben in \cite{sayood96} oder \cite{salomon2007}
\end{frame}

\subsection{Overlay}
\begin{frame}{\insertsection: \insertsubsection}
Bitte verwenden Sie Overlays nur moderat. Sie sind nicht immer hilfreich.
\begin{enumerate}
 \item<1> Eins 
 \item<2> Zwei
\end{enumerate}
\begin{itemize}
 \item Pause \pause
 \item Pause \pause
\end{itemize}
\begin{description}
 \item[dieses] \empho<3>{jenes}
 \item[solches] \empho<4>{und anderes}
\end{description}
Ende

Es gibt auch noch einige andere overlay Kommandos in der beamer Dokumentation.
\end{frame}

%%%%%%%%%%%%%%%%%%%%%%%%%%%%%%%%%%%%%%%%%%%%%%%%%%%%%%%%%%%%%%%%%%%%%%

\section{Zweiter Abschnitt} \label{sec:zweite}
\subsection{Overlay}
\begin{frame}{\insertsection: \insertsubsection}
 zweiter Abschnitt
\end{frame}



%%%%%%%%%%%%%%%%%%%%%%%%%%%%%%%%%%%%%%%%%%%%%%%%%%%%%%%%%%%%

\begin{frame}[label=finalSlide]
  \label{LastPage}%
  \begin{center}
    \vfill
    {\Large
    \textcolor{i6blue}{Danke für die Aufmerksamkeit!}
    }
%     \vfill
%     \inserttitle
    \vfill
    {\Large \insertauthor}
    \vfill
    \email{}
  \end{center}
\end{frame}

\end{document}
